% ============================================================
% C2 AI for Math — paper.tex
% 编译：pdflatex paper && bibtex paper && pdflatex paper && pdflatex paper
% 引用全部来自 references.bib（14 篇一手文献，作者名单已逐条核验）。
% 结构沿用老师大纲 9 节；论点：三层模型 + 语义规约层(Typed IR)。
% ============================================================
\documentclass[11pt]{article}

\usepackage[margin=1in]{geometry}
\usepackage{amsmath,amssymb,amsthm}
\usepackage{booktabs}
\usepackage{enumitem}
\usepackage{graphicx}
\usepackage[hidelinks]{hyperref}
\usepackage[round]{natbib}
\usepackage[T1]{fontenc}

\newcommand{\IR}{\ensuremath{\mathcal{IR}}}
\newcommand{\Nat}{\ensuremath{\mathbb{N}}}

\newtheorem{definition}{Definition}
\newtheorem{proposition}{Proposition}
\newtheorem{remark}{Remark}

\title{From Natural Language to Verifiable Reasoning:\\ A Semantic-Specification Framework for Reliable AI Mathematics}
\author{[Student Name]\\ \small AI for Math Challenge (C2)}
\date{\today}

\begin{document}
\maketitle

\begin{abstract}
Large language models (LLMs) can now solve a broad class of mathematical
problems, yet they remain unreliable: a plausible chain of reasoning can
conceal an invalid inference. We argue that reliability cannot be achieved by
post-hoc answer verification alone, nor by direct translation from natural
language into a formal proof assistant. Instead, we propose a three-layer
architecture that inserts a \emph{semantic-specification layer}---a typed
intermediate representation (\IR{})---between generation and verification. The
\IR{} carries a type system, an operational semantics, and composable
constraints, so that a mathematical claim is first ``compiled'' into a
checkable artifact before it reaches a theorem prover. We define a pipeline
(generation $\rightarrow$ semantic specification $\rightarrow$ verification)
and a bidirectional-consistency check that ties the natural-language statement,
the \IR{}, and the formal target to one another. We design a small but concrete
experiment on GSM8K and MATH that measures autoformalization and proof success
rates under this discipline. The central contribution is conceptual and
practical: moving the locus of trust from ``does the model sound right'' to
``does the \IR{} type-check and reduce to a proof obligation.''
\end{abstract}

\section{Introduction}
\label{sec:intro}

Large language models have made rapid progress on mathematical reasoning.
Chain-of-thought prompting elicits intermediate reasoning steps
\citep{wei2022chain}, self-consistency reduces variance by sampling many
reasoning paths \citep{wang2022selfconsistency}, and scaling has produced
systems such as Minerva that solve quantitative problems at a graduate level
\citep{lewkowycz2022solving}. On the public benchmarks GSM8K
\citep{cobbe2021training} and MATH \citep{hendrycks2021measuring}, strong
models now answer a large fraction of problems correctly, and
reinforcement-learning-based models such as DeepSeek-R1 demonstrate emergent
self-verification behaviours \citep{deepseek2025r1}.

These successes, however, measure \emph{answer correctness}, not
\emph{reasoning correctness}. A model that emits a confident but invalid step
still risks producing a wrong final answer, and on open-ended or unseen
mathematics there is no gold label to compare against. The core reliability
problem is therefore not ``can the model get the right number,'' but ``can we
trust that each step is sound.''

Two established remedies occupy opposite ends of a spectrum. On one end,
\emph{verifiers} are trained to score candidate solutions
\citep{cobbe2021training}, but they only rank plausibility and can be fooled by
correlated errors. On the other end, \emph{formal proof assistants} such as
Lean, Isabelle/HOL, and Coq provide machine-checkable certainty, but producing
a formal proof is expensive: automatic formalization of natural-language
mathematics remains hard and error-prone
\citep{wu2022autoformalization, azerbayev2023proofnet}. There is a wide,
under-explored gap between ``informal but fluent'' and ``formal but laborious.''

We propose to close part of this gap with a \emph{semantic-specification
layer}. Instead of jumping directly from a natural-language solution to a
formal proof, we first translate the solution into a \emph{typed intermediate
representation} (\IR{}) that makes the intended objects, operations, and
invariants explicit and mechanically checkable---but without the full burden of
a complete proof. This restates the reliability question in a tractable form:
an answer is trusted only if its \IR{} type-checks and its proof obligations
are dischargeable by a prover.

Our contributions are threefold. \textbf{(1)} We define a minimal but
well-typed \IR{} for elementary mathematical reasoning---a type system, an
operational semantics, and composable constraints---that is small enough to
reason about and rich enough to capture GSM8K/MATH-style problems.
\textbf{(2)} We specify a three-stage pipeline
(generation $\rightarrow$ semantic specification $\rightarrow$ verification)
together with a \emph{bidirectional-consistency} check that keeps the natural
language, the \IR{}, and the formal target mutually consistent. \textbf{(3)}
We design a concrete evaluation protocol on GSM8K and MATH that measures
autoformalization rate, type-check rate, and proof rate, following the
measurement conventions of ProofNet and LeanDojo
\citep{azerbayev2023proofnet, yang2023leandojo}.

The rest of the paper is organized as follows. Section~\ref{sec:model}
presents the three-layer model. Section~\ref{sec:bottleneck} analyses why
semantic specification is the critical bottleneck. Section~\ref{sec:method}
gives the \IR{} and the consistency checks. Section~\ref{sec:related} relates
our work to the literature. Section~\ref{sec:applications} sketches
applications. Section~\ref{sec:experiments} describes the experimental design,
Section~\ref{sec:discussion} discusses limitations, and
Section~\ref{sec:conclusion} concludes.
\section{The Three-Layer Model}
\label{sec:model}

We model reliable AI mathematics as a pipeline with three layers, each with a
distinct contract (Figure~\ref{fig:pipeline}).

\begin{description}[leftmargin=2.2em]
\item[Layer 1 -- Generation.] An LLM consumes a natural-language problem and
produces a candidate solution: a mixture of informal prose and computation.
This layer is fluent and general but has no soundness guarantee
\citep{wei2022chain, lewkowycz2022solving}.
\item[Layer 2 -- Semantic specification.] The candidate solution is compiled
into a typed intermediate representation that makes explicit: which objects
exist, which operations apply to them, what constraints hold, and which
sub-goals remain to be proven. This layer is the focus of the paper.
\item[Layer 3 -- Verification.] The \IR{}'s proof obligations are discharged
by a theorem prover or a symbolic checker, yielding a machine-checkable
certificate.
\end{description}

\begin{figure}[t]
\centering
\fbox{\begin{minipage}{0.9\linewidth}\centering
\textbf{Generation} (LLM, natural language)\\[2pt]
$\downarrow$ \emph{compile}\\[2pt]
\textbf{Semantic specification} (Typed \IR{}, checkable)\\[2pt]
$\downarrow$ \emph{reduce / discharge}\\[2pt]
\textbf{Verification} (Lean / Isabelle / Coq, machine-checked)
\end{minipage}}
\caption{The three-layer pipeline. Trust enters only at the semantic
specification layer, where the \IR{} is type-checked.}
\label{fig:pipeline}
\end{figure}

The key separation is \emph{where trust is established}. In the common
``LLM + verifier'' setup, trust is a probability estimated from a learned
scorer \citep{cobbe2021training}. In the ``LLM + prover'' setup, trust is a
proof, but the model must produce the proof in one step
\citep{wu2022autoformalization}. Our model inserts a third, intermediate locus
of trust: a typed artifact that can be checked mechanically and independently
of how it was generated. This does not remove the difficulty---writing the
compiler from natural language to \IR{} is itself hard---but it makes the
difficulty explicit and checkable.

\section{Why Semantic Specification Is the Bottleneck}
\label{sec:bottleneck}

We claim that the semantic-specification layer is the \emph{critical} part of
the pipeline, for three reasons.

First, \emph{the hard part of mathematics is not the final answer but the
mapping from informal intent to formal object}. A word problem such as ``how
many apples remain'' requires deciding that apples are natural numbers, that
subtraction is the right operation, and that the result is non-negative. This
mapping is precisely what natural language leaves implicit and what a proof
assistant requires explicit. Errors here---wrong type, wrong operation, missing
precondition---are the dominant failure mode of direct autoformalization
\citep{wu2022autoformalization}.

Second, \emph{post-hoc verification cannot recover from an upstream
mis-specification}. If the model has decided to use division where
multiplication was intended, a verifier that only scores the final answer may
not detect the error, and may even reinforce it. Only a representation that
fixes the intended semantics can expose the mismatch.

Third, \emph{full formal proof is overkill for many downstream uses}. In
education and rapid problem solving we often need \emph{enough} rigor to catch
invalid steps, not a complete machine-checked proof of a difficult theorem.
The \IR{} is designed to capture exactly this intermediate level: checkable
rigor without full proof burden.
\section{Method: A Typed Intermediate Representation}
\label{sec:method}

We now make the semantic-specification layer concrete. The design goal is
\emph{minimality with expressiveness}: enough structure to catch invalid
inferences on GSM8K/MATH-style problems, but simple enough to reason about and
implement.

\subsection{The \IR{}}
\label{sec:ir}

The \IR{} is a small typed term language over a set of sorts
$\{\Nat, \mathbb{R}, \texttt{Bool}\}$ (extensible). Its syntax is an abstract
syntax tree (AST) built from variables, literals, and typed operations.

\begin{definition}[Typed term]
A well-typed \IR{} term is an AST node that carries an explicit sort and whose
children are themselves well-typed. Every operation has a type signature;
applying an operation to arguments of the wrong sort is a \emph{type error} and
is rejected before verification.
\end{definition}

\begin{definition}[Operational semantics]
Each primitive operation has a small-step reduction rule. A term is
\emph{reducible} if some subterm matches a reduction rule; a term is in
\emph{normal form} if no rule applies and no proof obligation remains open.
\end{definition}

The crucial design element is that the \IR{} is \emph{composable}: a solution
is a tree of smaller, independently checkable facts. This mirrors how human
mathematicians build proofs, and it lets the checker localise a failure to the
exact subterm that mis-typed.

\begin{proposition}[Localisation]
If an \IR{} term fails to type-check, the failing subterm can be identified in
time linear in the size of the AST.
\end{proposition}

\subsection{From Generation to Specification}
\label{sec:compile}

The compiler from Layer~1 to Layer~2 is an LLM prompt-and-check loop, not a
hand-written parser: the model proposes a typed AST from the natural-language
solution, and a deterministic checker rejects it if it fails to type-check or
violates an explicit constraint. This is the same ``propose then verify''
discipline used in neural theorem provers
\citep{wang2023legoprover, yang2023leandojo}, but applied one level earlier, at
specification rather than proof.

\subsection{Bidirectional Consistency}
\label{sec:consistency}

Because the \IR{} is generated, it can drift from the natural-language
statement it is supposed to represent. We therefore require a
\emph{bidirectional-consistency check} with two directions.

\begin{itemize}
\item \textbf{Forward check} (natural language $\rightarrow$ \IR{}): a model
judges, for each \IR{} leaf and operation, whether it faithfully captures the
corresponding phrase of the original problem. Disagreements are flagged as
\emph{specification drift}.
\item \textbf{Backward check} (\IR{} $\rightarrow$ formal target): the \IR{}'s
proof obligations are mapped to a theorem-prover goal, and the mapping must be
type-preserving. This check is mechanical once the target calculus is fixed.
\end{itemize}

Together these checks tie the three layers together: the natural-language
statement, the \IR{}, and the formal target cannot silently diverge.
\section{Related Work}
\label{sec:related}

Our work sits between two active literatures.

\textbf{LLM mathematical reasoning (generation side).} Chain-of-thought
prompting \citep{wei2022chain} and self-consistency
\citep{wang2022selfconsistency} improve accuracy by eliciting and aggregating
reasoning paths, but they provide no soundness guarantee. Plan-and-Solve
prompting \citep{wang2023plan} adds explicit planning, which is conceptually
close to our compilation step, but its plan is informal prose rather than a
typed artifact. Verifier-based training \citep{cobbe2021training} and scaled
systems such as Minerva \citep{lewkowycz2022solving} and DeepSeek-R1
\citep{deepseek2025r1} raise accuracy, yet they remain probabilistic: none of
them produces a checkable intermediate representation of the reasoning itself.

\textbf{Formalisation and theorem proving (verification side).}
Autoformalization \citep{wu2022autoformalization} and ProofNet
\citep{azerbayev2023proofnet} study translating natural-language mathematics
into Lean, exposing precisely the ``semantic gap'' we address, but they target
full formal proofs directly. LeanDojo \citep{yang2023leandojo} and LEGO-Prover
\citep{wang2023legoprover} use retrieval and growing tactic libraries to make
neural theorem proving practical, again at the proof level. On the
verification-engineering side, Isabelle/HOL provides a mature framework for
reasoning about programs and specifications \citep{schirmer2006verification},
KeYmaera X demonstrates tactic-based proof for hybrid systems
\citep{fulton2015keymaera}, and Coq's self-verification of its own type checker
\citep{sozeau2019coqcoqcorrect} shows that even the checking machinery can be
made trustworthy.

Our contribution is to occupy the middle layer that these works leave open:
a typed, checkable specification that is cheaper than a full proof but more
rigorous than informal reasoning, with explicit bidirectional consistency
between the natural-language input, the \IR{}, and the formal target.

\section{Applications}
\label{sec:applications}

The three-layer model is not tied to any single prover or domain. We sketch
three applications.

\begin{itemize}
\item \textbf{Automated theorem proving.} The \IR{} serves as a
\emph{specification first, proof second} front end: a system formalises a
problem into \IR{}, checks the specification, and only then invokes the prover
on the reduced obligations, reducing wasted proof search
\citep{yang2023leandojo}.
\item \textbf{Problem solving.} For word problems on GSM8K/MATH, the \IR{}
acts as a ``show your work'' layer: a solution is accepted only if it
type-checks, giving a lightweight correctness signal without a full formal
proof \citep{cobbe2021training, hendrycks2021measuring}.
\item \textbf{AI-assisted education.} Because the \IR{} localises failures to
specific subterms, it can produce targeted feedback---``this step divided by a
quantity that could be zero''---which is more instructive than a binary
right/wrong verdict.
\end{itemize}
\section{Experimental Design}
\label{sec:experiments}

We design a small, reproducible evaluation that isolates the effect of the
semantic-specification layer. Following ProofNet and LeanDojo conventions
\citep{azerbayev2023proofnet, yang2023leandojo}, we report three rates on a
sample of GSM8K and MATH problems
\citep{cobbe2021training, hendrycks2021measuring}.

\begin{description}[leftmargin=2.2em]
\item[Autoformalization rate.] The fraction of problems for which the model
produces a well-typed \IR{} from the natural-language statement
\citep{wu2022autoformalization}.
\item[Type-check rate.] The fraction of proposed \IR{} terms that pass the
deterministic checker (Section~\ref{sec:ir}).
\item[Proof rate.] The fraction of proof obligations discharged by the target
prover (Lean), the metric used by neural theorem provers
\citep{wang2023legoprover, yang2023leandojo}.
\end{description}

We also measure the \emph{consistency} between the final answer accepted by the
\IR{} and the ground-truth answer, and we record the number of rejected
proposals per problem as a cost signal. The primary comparison is
\emph{ablation}: the same model with and without the \IR{} type-check gate, to
test the hypothesis that the gate raises reasoning correctness without
requiring a full proof.

\section{Discussion}
\label{sec:discussion}

\textbf{Scope.} Our \IR{} currently covers arithmetic and elementary algebra
sorts; richer domains (analysis, geometry, hybrid systems
\citep{fulton2015keymaera}) require extending the sort system and operational
semantics, which we treat as future work.

\textbf{Limitations.} The compiler from natural language to \IR{} is itself an
LLM and can mis-specify; the bidirectional-consistency check mitigates but does
not eliminate this risk. Type-checking guarantees that a term is
\emph{well-typed}, not that it is \emph{the intended} term. Full assurance
still requires the proof layer \citep{sozeau2019coqcoqcorrect}.

\textbf{Ethics.} We do not claim that the framework makes AI mathematics
infallible. It lowers, but does not remove, the risk of confident incorrect
answers; in high-stakes settings a human review of the \IR{} remains advisable.

\section{Conclusion}
\label{sec:conclusion}

We argued that reliable AI mathematics needs a checkable intermediate layer
between fluent generation and costly formal proof, and we made that layer
concrete as a typed intermediate representation with an operational semantics,
composable constraints, and bidirectional consistency checks. The framework
recasts the reliability question from ``does the model sound right'' to ``does
the specification type-check and reduce to a proof obligation,'' which is
mechanically decidable one level before full proof. Our evaluation protocol on
GSM8K and MATH is designed to measure whether this discipline measurably
improves reasoning correctness at lower cost than complete formalization. We
believe the semantic-specification layer is a necessary, and currently
under-investigated, building block for trustworthy AI mathematics.

\bibliographystyle{plainnat}
\bibliography{references}

\end{document}

